% Generated by ddl_bna.py
\textsc{Agencia} & \texttt{(id\_\allowbreak{}agencia)} & \texttt{distrito} no determina \texttt{ciudad}: hay distritos homónimos en distintas ciudades (Miraflores en Lima y en Arequipa). \\
\textsc{Cuenta\_Contable} & \texttt{(codigo\_\allowbreak{}contable)} & Se retiró \texttt{nivel} en 3FN; \texttt{naturaleza} no depende de \texttt{codigo\_padre}, porque existen cuentas correctoras de naturaleza opuesta a la de su cuenta padre. \\
\textsc{Asiento\_Contable} & \texttt{(id\_\allowbreak{}asiento)}; \texttt{(numero\_\allowbreak{}asiento)}; \textit{únicas con nulos:} \texttt{(id\_\allowbreak{}asiento\_\allowbreak{}revertido)} & \texttt{numero\_asiento} es clave candidata; \texttt{id\_asiento\_revertido} es única pero admite nulos, por lo que no es clave y no determina nada. \texttt{estado\_asiento} se retiró por derivarse de otra fila. \\
\textsc{Cliente} & \texttt{(cod\_\allowbreak{}cliente)}; \texttt{(tipo\_\allowbreak{}documento, numero\_\allowbreak{}documento)}; \texttt{(correo)}; \texttt{(celular)}; \texttt{(id\_\allowbreak{}verificacion\_\allowbreak{}captacion)} & Cuatro claves candidatas más (documento, correo, celular y verificación); todo determinante es superclave. \\
\textsc{Biometria\_Cliente} & \texttt{(id\_\allowbreak{}biometria)} & Clave simple; ningún atributo no primo determina a otro. \\
\textsc{Codigo\_Verificacion} & \texttt{(id\_\allowbreak{}codigo)} & \texttt{estado\_codigo} no se deriva de \texttt{intentos\_fallidos} ni de \texttt{fecha\_expiracion}: VALIDADO, PENDIENTE y EXPIRADO dependen del momento y de la respuesta del cliente. \\
\textsc{Validacion\_Identidad} & \texttt{(id\_\allowbreak{}validacion)} & Clave simple; ningún atributo no primo determina a otro. \\
\textsc{Cuenta\_Ahorro} & \texttt{(num\_\allowbreak{}cuenta)}; \texttt{(cci)} & \texttt{cci} es clave candidata; aunque codifica la oficina y el número de cuenta, todo lo que determina lo determina como superclave. \\
\textsc{Movimiento} & \texttt{(uuid\_\allowbreak{}transaccion)}; \texttt{(id\_\allowbreak{}asiento)} & \texttt{id\_asiento} es clave candidata (relación 1:1 con el asiento). \\
\textsc{Canal} & \texttt{(id\_\allowbreak{}canal)}; \texttt{(nombre)} & Clave simple; ningún atributo no primo determina a otro. \\
\textsc{Tarjeta} & \texttt{(id\_\allowbreak{}tarjeta)}; \texttt{(pan\_\allowbreak{}token)} & Tras retirar \texttt{tipo\_tarjeta} y \texttt{vencimiento\_mm\_aa}, ningún atributo no primo determina a otro. \\
\textsc{Matriz\_Token} & \texttt{(id\_\allowbreak{}matriz)}; \texttt{(numero\_\allowbreak{}serie)} & Clave simple; ningún atributo no primo determina a otro. \\
\textsc{Token\_Pago} & \texttt{(id\_\allowbreak{}token\_\allowbreak{}pago)}; \texttt{(valor\_\allowbreak{}token)} & Clave simple; ningún atributo no primo determina a otro. \\
\textsc{Intento\_Autenticacion} & \texttt{(id\_\allowbreak{}intento)} & \texttt{id\_canal} solo tiene valor cuando no hay terminal (restricción de arco exclusivo), de modo que no depende de \texttt{id\_terminal}. \\
\textsc{Autorizacion} & \texttt{(uuid\_\allowbreak{}autorizacion)}; \textit{únicas con nulos:} \texttt{(uuid\_\allowbreak{}transaccion)} & Igual que el intento respecto de terminal y canal; \texttt{codigo\_respuesta} no determina el estado (00 corresponde a AUTORIZADA y a DISPENSADA). \texttt{uuid\_transaccion} es única pero admite nulos (las consultas no generan movimiento). \\
\textsc{Historial\_Estado\_Tarjeta} & \texttt{(id\_\allowbreak{}historial)} & Clave simple; ningún atributo no primo determina a otro. \\
\textsc{Evento\_Salida} & \texttt{(id\_\allowbreak{}evento)} & Clave simple; ningún atributo no primo determina a otro. \\
\textsc{Solicitud\_Credito} & \texttt{(id\_\allowbreak{}solicitud)} & Clave simple; ningún atributo no primo determina a otro. \\
\textsc{Historial\_Estado\_Solicitud} & \texttt{(id\_\allowbreak{}historial)} & Clave simple; ningún atributo no primo determina a otro. \\
\textsc{Garantia} & \texttt{(id\_\allowbreak{}garantia)} & Clave simple; ningún atributo no primo determina a otro. \\
\textsc{Fiador} & \texttt{(id\_\allowbreak{}fiador)}; \texttt{(id\_\allowbreak{}persona\_\allowbreak{}externa)}; \texttt{(tipo\_\allowbreak{}documento, numero\_\allowbreak{}documento)} & Clave simple; ningún atributo no primo determina a otro. \\
\textsc{Calendario\_Pago} & \texttt{(id\_\allowbreak{}calendario)}; \texttt{(id\_\allowbreak{}solicitud, version)} & (\texttt{id\_solicitud}, \texttt{version}) es clave candidata. \\
\textsc{Situacion\_Financiera} & \texttt{(id\_\allowbreak{}consulta)} & Clave simple; ningún atributo no primo determina a otro. \\
\textsc{Evaluacion} & \texttt{(id\_\allowbreak{}evaluacion)}; \texttt{(id\_\allowbreak{}solicitud, ronda)} & (\texttt{id\_solicitud}, \texttt{ronda}) es clave candidata. \texttt{score\_calculado} y \texttt{limite\_aprobado} se registran como resultado del motor en esa ronda: el modelo de score es externo y versionado, así que \texttt{id\_consulta} no los determina. \\
\textsc{Regla\_Evaluacion} & \texttt{(id\_\allowbreak{}regla)}; \texttt{(codigo\_\allowbreak{}rn)}; \textit{únicas con nulos:} \texttt{(ratio\_\allowbreak{}min, ratio\_\allowbreak{}max)} & \texttt{codigo\_rn} es clave candidata; el tramo (\texttt{ratio\_min}, \texttt{ratio\_max}) es único, pero un extremo puede ser nulo (tramos abiertos), por lo que no es clave. \texttt{accion} depende de la regla, no del tramo por separado. \\
\textsc{Analista} & \texttt{(id\_\allowbreak{}analista)}; \texttt{(codigo\_\allowbreak{}empleado)} & Clave simple; ningún atributo no primo determina a otro. \\
\textsc{Decision\_Analista} & \texttt{(id\_\allowbreak{}decision)}; \texttt{(id\_\allowbreak{}evaluacion)} & \texttt{id\_evaluacion} es clave candidata (relación 1:1). \\
\textsc{Orden\_Desembolso} & \texttt{(id\_\allowbreak{}orden)}; \texttt{(codigo\_\allowbreak{}orden)}; \texttt{(idempotency\_\allowbreak{}key)}; \textit{únicas con nulos:} \texttt{(uuid\_\allowbreak{}transaccion\_\allowbreak{}abono)} & \texttt{codigo\_orden} e \texttt{idempotency\_key} son claves candidatas; \texttt{request\_hash} depende de la petición, identificada por la clave de idempotencia. \texttt{uuid\_transaccion\_abono} es única pero admite nulos (solo la orden confirmada abona). \\
\textsc{Validacion\_Desembolso} & \texttt{(id\_\allowbreak{}validacion)} & Clave simple; ningún atributo no primo determina a otro. \\
\textsc{Evento\_Entrega} & \texttt{(id\_\allowbreak{}evento)} & \texttt{version\_esquema} no depende de \texttt{tipo\_evento}: los eventos antiguos conservan la versión del contrato con la que se enviaron. \\
\textsc{Solicitud\_Reversion} & \texttt{(id\_\allowbreak{}reversion)}; \texttt{(idempotency\_\allowbreak{}key)}; \textit{únicas con nulos:} \texttt{(uuid\_\allowbreak{}transaccion\_\allowbreak{}extorno)} & \texttt{idempotency\_key} es clave candidata; \texttt{uuid\_transaccion\_extorno} es única pero admite nulos (solo la reversión confirmada genera extorno). \\
\textsc{Credito} & \texttt{(cod\_\allowbreak{}credito)}; \texttt{(id\_\allowbreak{}orden)}; \texttt{(id\_\allowbreak{}calendario)} & \texttt{id\_orden} e \texttt{id\_calendario} son claves candidatas (relaciones 1:1). \\
\textsc{Pago} & \texttt{(id\_\allowbreak{}pago)}; \texttt{(referencia\_\allowbreak{}transaccion)} & Clave simple; ningún atributo no primo determina a otro. \\
\textsc{Mora} & \texttt{(id\_\allowbreak{}mora)}; \texttt{(id\_\allowbreak{}calendario, numero\_\allowbreak{}cuota, fecha\_\allowbreak{}calculo)} & (\texttt{id\_calendario}, \texttt{numero\_cuota}, \texttt{fecha\_calculo}) es clave candidata; \texttt{dias\_atraso} se retiró en 3FN. \\
\textsc{Evento\_Auditoria} & \texttt{(id\_\allowbreak{}evento)} & El tipo del registro auditado se obtiene de ENTIDAD\_AUDITADA; \texttt{numero\_cuota} solo tiene valor cuando el registro es una cuota y no depende de \texttt{id\_entidad}, porque un calendario tiene varias cuotas. \\
\textsc{Detalle\_Asiento} & \texttt{(id\_\allowbreak{}asiento, nro\_\allowbreak{}linea)} & Clave compuesta; los atributos no primos dependen de toda la clave y no entre sí. \\
\textsc{Consumo\_Diario} & \texttt{(num\_\allowbreak{}cuenta, id\_\allowbreak{}canal, fecha\_\allowbreak{}corte)} & Clave compuesta; los atributos no primos dependen de toda la clave y no entre sí. \\
\textsc{Limite\_Canal} & \texttt{(id\_\allowbreak{}canal, moneda)} & Tras trasladar la regla a CANAL, el tope y la habilitación dependen del par canal--moneda. \\
\textsc{Credencial} & \texttt{(id\_\allowbreak{}tarjeta, tipo\_\allowbreak{}credencial)} & Tras extraer \texttt{longitud}, los atributos dependen del par tarjeta--tipo. \\
\textsc{Matriz\_Coordenada} & \texttt{(id\_\allowbreak{}matriz, fila, columna)} & Clave compuesta; los atributos no primos dependen de toda la clave y no entre sí. \\
\textsc{Cuota} & \texttt{(id\_\allowbreak{}calendario, numero\_\allowbreak{}cuota)} & \texttt{estado\_cuota} depende de la fecha de consulta además de los saldos (VENCIDA), por lo que no hay dependencia funcional entre atributos. \\
\textsc{Titularidad} & \texttt{(cod\_\allowbreak{}cliente, num\_\allowbreak{}cuenta)} & Sin atributos no primos dependientes de parte de la clave tras retirar \texttt{rol\_titular} en 2FN. \\
\textsc{Tarjeta\_Cuenta} & \texttt{(id\_\allowbreak{}tarjeta, num\_\allowbreak{}cuenta)} & Clave compuesta; los atributos no primos dependen de toda la clave y no entre sí. \\
\textsc{Fiador\_Solicitud} & \texttt{(id\_\allowbreak{}fiador, id\_\allowbreak{}solicitud)}; \texttt{(id\_\allowbreak{}verificacion)} & Clave compuesta; los atributos no primos dependen de toda la clave y no entre sí. \\
\textsc{Aplicacion\_Pago} & \texttt{(id\_\allowbreak{}pago, id\_\allowbreak{}calendario, numero\_\allowbreak{}cuota)} & Clave compuesta; los atributos no primos dependen de toda la clave y no entre sí. \\
\textsc{Canal\_Validacion} & \texttt{(id\_\allowbreak{}canal, tipo\_\allowbreak{}validacion)} & Clave compuesta; los atributos no primos dependen de toda la clave y no entre sí. \\
\textsc{Tipo\_Credencial} & \texttt{(tipo\_\allowbreak{}credencial)} & Clave simple; ningún atributo no primo determina a otro. \\
\textsc{Verificacion\_Captacion} & \texttt{(id\_\allowbreak{}verificacion\_\allowbreak{}captacion)} & Clave simple; ningún atributo no primo determina a otro. \\
\textsc{Tipo\_Operacion} & \texttt{(tipo\_\allowbreak{}operacion)} & Clave simple; ningún atributo no primo determina a otro. \\
\textsc{Bin\_Tarjeta} & \texttt{(bin)} & Clave simple; ningún atributo no primo determina a otro. \\
\textsc{Terminal} & \texttt{(id\_\allowbreak{}terminal)} & \texttt{id\_agencia} no determina \texttt{id\_canal}: una agencia tiene ventanillas, agentes y cajeros. \\
\textsc{Tipo\_Validacion\_Desembolso} & \texttt{(tipo\_\allowbreak{}validacion)} & Clave simple; ningún atributo no primo determina a otro. \\
\textsc{Tipo\_Evento\_Salida} & \texttt{(tipo\_\allowbreak{}evento)} & Clave simple; ningún atributo no primo determina a otro. \\
\textsc{Entidad\_Auditada} & \texttt{(id\_\allowbreak{}entidad)} & Un único atributo no primo: el tipo del registro auditado. \\
\textsc{Moneda} & \texttt{(cod\_\allowbreak{}moneda)} & Catálogo: código como clave y atributos que describen solo a ese código. \\
\textsc{Tipo\_Documento} & \texttt{(cod\_\allowbreak{}tipo\_\allowbreak{}documento)} & Catálogo: código como clave y atributos que describen solo a ese código. \\
\textsc{Tipo\_Biometria} & \texttt{(cod\_\allowbreak{}tipo\_\allowbreak{}biometria)} & Catálogo: código como clave y atributos que describen solo a ese código. \\
\textsc{Tipo\_Contacto} & \texttt{(cod\_\allowbreak{}tipo\_\allowbreak{}contacto)} & Catálogo: código como clave y atributos que describen solo a ese código. \\
\textsc{Tipo\_Validacion} & \texttt{(cod\_\allowbreak{}tipo\_\allowbreak{}validacion)}; \texttt{(endpoint)} & Catálogo: código como clave y atributos que describen solo a ese código. \\
\textsc{Regimen\_Cuenta} & \texttt{(cod\_\allowbreak{}regimen)} & \texttt{firma\_conjunta} y \texttt{min\_titulares} no se determinan entre sí (la indistinta y la conjunta exigen dos titulares). \\
\textsc{Estado\_Cuenta} & \texttt{(cod\_\allowbreak{}estado)} & Catálogo: código como clave y atributos que describen solo a ese código. \\
\textsc{Estado\_Tarjeta} & \texttt{(cod\_\allowbreak{}estado)} & Catálogo: código como clave y atributos que describen solo a ese código. \\
\textsc{Estado\_Cliente} & \texttt{(cod\_\allowbreak{}estado)} & Catálogo: código como clave y atributos que describen solo a ese código. \\
\textsc{Estado\_Solicitud} & \texttt{(cod\_\allowbreak{}estado)} & Catálogo: código como clave y atributos que describen solo a ese código. \\
\textsc{Estado\_Orden} & \texttt{(cod\_\allowbreak{}estado)} & Catálogo: código como clave y atributos que describen solo a ese código. \\
\textsc{Estado\_Credito} & \texttt{(cod\_\allowbreak{}estado)} & Catálogo: código como clave y atributos que describen solo a ese código. \\
\textsc{Estado\_Cuota} & \texttt{(cod\_\allowbreak{}estado)} & Catálogo: código como clave y atributos que describen solo a ese código. \\
\textsc{Codigo\_Respuesta} & \texttt{(cod\_\allowbreak{}respuesta)} & Catálogo: código como clave y atributos que describen solo a ese código. \\
\textsc{Origen\_Cambio} & \texttt{(cod\_\allowbreak{}origen)} & Catálogo: código como clave y atributos que describen solo a ese código. \\
\textsc{Tipo\_Garantia} & \texttt{(cod\_\allowbreak{}tipo\_\allowbreak{}garantia)} & Catálogo: código como clave y atributos que describen solo a ese código. \\
\textsc{Metodo\_Amortizacion} & \texttt{(cod\_\allowbreak{}metodo)} & Catálogo: código como clave y atributos que describen solo a ese código. \\
\textsc{Resultado\_Crediticio} & \texttt{(cod\_\allowbreak{}resultado)} & Catálogo: código como clave y atributos que describen solo a ese código. \\
\textsc{Modulo} & \texttt{(id\_\allowbreak{}modulo)}; \texttt{(nombre)} & Catálogo: código como clave y atributos que describen solo a ese código. \\
\textsc{Pasarela\_Pago} & \texttt{(id\_\allowbreak{}pasarela)} & Catálogo: código como clave y atributos que describen solo a ese código. \\
